% NEW ACRONYM — add to formatting/Acronyms.tex:
% \DeclareAcronym{RTX}{short=RTX, long={NVIDIA real-time ray tracing}, tag=abbrev, sort=RTX}

%===================================================================%
\section{NVIDIA Isaac Sim and IsaacLab Architecture}
\label{sec:isaac_sim_isaaclab_architecture}
%===================================================================%

NVIDIA Isaac Sim (version~5.1) is employed as the primary simulation environment for the learning-based manipulation pipeline considered in this thesis~\cite{isaac_sim_510_overview,isaac_sim_510_release_notes}.
Isaac Sim is provided as a reference application built on NVIDIA Omniverse for the development, simulation, and testing of AI-driven robotic systems in physically-based virtual environments~\cite{isaac_sim_510_overview}.
Its architecture combines a \ac{USD}-centric scene representation, a high-fidelity \emph{GPU}-based physics backend, and an RTX-enabled rendering pipeline to support both control-centric simulation and perception-oriented data generation~\cite{isaac_sim_510_overview}.
In addition, robotics interfaces are exposed through dedicated extensions and bridge APIs, including integration paths to \ac{ROS}~2 for interoperability with established robotics middleware ecosystems~\cite{isaac_sim_510_overview,isaac_sim_ros2_reference_architecture_510}.

\subsection{Isaac Sim as an Omniverse-based simulation platform}
Isaac Sim is built around Omniverse concepts of a composable application framework and an extensible scene graph that is represented and manipulated through \ac{USD}~\cite{isaac_sim_510_overview}.
A key architectural implication is that simulation assets, sensors, and tools are integrated as modular extensions, enabling application functionality to be assembled and updated without altering the underlying scene description~\cite{isaac_sim_510_overview}.
In Isaac Sim, interoperability with common robotics asset formats is supported through import workflows for \ac{URDF} and \ac{MJCF}, which are mapped into a unified \ac{USD} representation and can subsequently be tuned within a consistent authoring and simulation workflow~\cite{isaac_sim_510_overview}.
This unification is particularly relevant for research pipelines where heterogeneous assets (robots, fixtures, sensors, and task-specific objects) must be combined into a single simulation stage while preserving a reproducible configuration state~\cite{isaac_sim_510_overview}.

The simulation stack is exposed through multiple interaction modalities.
On the one hand, interactive authoring and debugging are supported through Omniverse-centric tools such as Omnigraph-based orchestration and application-level profiling utilities~\cite{isaac_sim_510_overview}.
On the other hand, the same scene can be executed in a headless manner without rendering, which is essential for high-throughput learning workloads where visual fidelity is not required for every training step~\cite{isaac_sim_510_physics_fundamentals}.
The availability of both interactive and headless operation modes enables a practical separation between (i) rapid iteration on scene and asset correctness and (ii) large-scale policy training and evaluation~\cite{isaac_sim_510_physics_fundamentals}.

\subsection{Core simulation and rendering components}
\subsubsection{Physics representation and PhysX backend}
The physics interface in Isaac Sim is defined through a combination of \ac{USD} Physics schemas and PhysX-specific schemas that annotate scene primitives with physical properties and simulation semantics~\cite{isaac_sim_510_physics_fundamentals}.
During simulation, these schema annotations are interpreted to instantiate corresponding simulation objects in the PhysX backend, thereby decoupling the authoring representation from the runtime solver state~\cite{isaac_sim_510_physics_fundamentals}.
Within this architecture, rigid bodies, collisions, joints, and articulated mechanisms are expressed at the \ac{USD} level while being numerically integrated by PhysX during stepping~\cite{isaac_sim_510_physics_fundamentals}.
Consequently, a consistent representation is maintained across interactive editing, scripted environment generation, and batched training execution, provided that schema-consistent authoring is enforced~\cite{isaac_sim_510_physics_fundamentals}.

PhysX is explicitly designed to support \ac{GPU}-accelerated simulation pathways.
In the PhysX \ac{GPU} simulation pipeline, collision detection and rigid-body dynamics can be offloaded to the \ac{GPU} to increase throughput for scenes that are compatible with the corresponding feature subset~\cite{physx_541_gpu_rigid_bodies}.
In addition, \ac{GPU}-centric simulation features include particle- and deformable-related capabilities such as soft bodies and particles, with the documentation explicitly noting that these features are not available in CPU-only mode~\cite{physx_541_gpu_rigid_bodies}.
For learning workloads, the resulting throughput advantage is typically realized by increasing the number of parallel environments rather than increasing the physical fidelity of a single environment, which aligns with the data-hungry nature of policy optimization in \acp{MDP}~\cite{Makoviychuk2021IsaacGym}.

\subsubsection{Photorealistic rendering with RTX}
Beyond physics, Isaac Sim provides multi-sensor rendering capabilities that are designed to scale to industrially sized scenes through direct \ac{GPU} access~\cite{isaac_sim_510_overview}.
For perception-oriented tasks, the Omniverse rendering stack offers RTX-based modes that include real-time rendering and interactive path tracing variants, as well as dedicated debugging views and render settings tailored to performance analysis~\cite{omniverse_rtx_renderer_doc}.
In Isaac Sim, this rendering capability is integrated with simulated sensors such as cameras and LiDARs, enabling the generation of perception streams that are synchronized with the physics state and can be employed for observation construction, validation of scene realism, or synthetic data generation~\cite{isaac_sim_510_overview}.
When rendering is not required (e.g., state-based control policies), the simulation can be executed without rendering to maximize training throughput~\cite{isaac_sim_510_physics_fundamentals}.

\subsubsection{ROS 2 integration and middleware bridging}
To interface with external robotics software stacks, Isaac Sim provides \ac{ROS}~2 integration via a bridge-based architecture.
In the recommended setup, \ac{ROS}~2 nodes are executed in a separate process (typically a container), while Isaac Sim runs its own \ac{ROS}~2 bridge extension that publishes and subscribes to \ac{ROS}~2 topics according to a defined data flow~\cite{isaac_sim_ros2_reference_architecture_510}.
This process-level separation supports robust integration with external tooling (e.g., logging, visualization, and hardware interfaces) while isolating simulation execution and preserving reproducible simulation states~\cite{isaac_sim_ros2_reference_architecture_510}.
The bridge-based approach is particularly relevant for sim-to-real transfer pipelines where the same \ac{ROS}~2 interfaces can be reused across simulation and real systems with minimal adaptation~\cite{isaac_sim_ros2_reference_architecture_510}.

\subsubsection{Summary of Isaac Sim architectural layers}
The principal layers relevant to the present thesis are summarized in \autoref{fig:isaac_stack_overview}.
A layered view is helpful because learning-centric simulation involves repeated traversal of the loop \emph{(state acquisition) $\rightarrow$ (policy inference) $\rightarrow$ (action application) $\rightarrow$ (physics stepping)}, while rendering and middleware integration remain optional subsystems depending on the observation design and deployment strategy~\cite{isaac_sim_510_physics_fundamentals,isaac_sim_ros2_reference_architecture_510}.

\begin{figure}[ht!]
\centering
\fbox{\parbox{0.8\textwidth}{\centering \vspace{0.1in}Placeholder for \textit{layered architecture of Isaac Sim and IsaacLab} image\vspace{0.1in}}}
\caption[Layered view of Isaac Sim and IsaacLab]{Layered view of the simulation and learning stack used in this thesis. The scene and asset layer is represented in \ac{USD}. Physics semantics are expressed via \ac{USD} Physics and PhysX schemas and executed by the PhysX backend. Rendering is provided through RTX-enabled Omniverse rendering modes. IsaacLab builds on Isaac Sim to provide reusable task abstractions and scalable training workflows.}\label{fig:isaac_stack_overview}
\end{figure}

\subsection{IsaacLab as a unified robotics learning framework}
IsaacLab is provided as a robotics learning framework built upon Isaac Sim and is intended to consolidate task design workflows and training utilities for robotics research~\cite{isaaclab_overview_doc,isaac_sim_510_overview}.
Within IsaacLab, task implementations and training scripts are organized such that simulation components (robots, sensors, assets, and physics settings) are described through configurable modules that can be reused and composed across tasks~\cite{isaaclab_overview_doc}.
The documentation explicitly positions IsaacLab as a successor to earlier Omniverse-based research tooling and notes that its development was initiated from the Orbit framework~\cite{isaaclab_overview_doc}.
Orbit itself introduces a modular simulation and learning framework on Isaac Sim, targeting reproducible and scalable robot learning through composable environment definitions and standardized training interfaces~\cite{Mittal2023Orbit}.

From an architectural perspective, IsaacLab can be interpreted as a consolidation of complementary design goals previously addressed by distinct NVIDIA research stacks.
Isaac Gym demonstrates that high-throughput robot learning can be achieved through \ac{GPU}-accelerated physics simulation and massive parallelization, thereby enabling the collection of large batches of experience for deep \ac{RL}~\cite{Makoviychuk2021IsaacGym}.
Orbit demonstrates that environment design on top of Isaac Sim benefits from modular abstractions that separate scene composition, command generation, observation construction, and reward specification to improve reusability and experimental control~\cite{Mittal2023Orbit}.
In IsaacLab, these goals are pursued within a single framework that is tightly integrated with Isaac Sim’s \ac{USD} and extension ecosystem, while also providing learning-oriented utilities such as task templates and scalable training workflows~\cite{isaaclab_overview_doc}. % TODO: find citation explicitly stating "combines Isaac Gym and Orbit"

\subsection{IsaacLab task design as an MDP workflow}
A central design principle in IsaacLab is that a task is formalized as a \ac{MDP}, where state transitions are induced by applying actions in a simulator and collecting observations and rewards.
To structure this workflow, IsaacLab provides manager abstractions that assign responsibility for individual \ac{MDP} components (e.g., observation construction, action application, reward computation, and termination logic) to dedicated modules~\cite{isaaclab_mdp_api}.
In particular, the documentation describes that different managers are responsible for observations, actions, commands, events, rewards, terminations, and curriculum scheduling~\cite{isaaclab_mdp_api}.
This separation supports reuse because the same robot and scene configuration can be paired with different observation or reward definitions, and it supports experimental reproducibility because changes can be localized to a specific manager configuration rather than being distributed across monolithic environment code~\cite{isaaclab_mdp_api}.

Within the resulting workflow, the simulation state is advanced in discrete time steps.
At each step, actions produced by a policy are applied to the simulator (e.g., as target joint positions, velocities, or torques), the physics is stepped, and the next observation and reward are computed from the updated simulator state~\cite{isaaclab_mdp_api}.
Because the \ac{MDP} definition is encapsulated in configuration and manager modules, training pipelines can be instantiated programmatically and scaled over a large number of parallel environments, which is a prerequisite for data-efficient policy learning with modern deep \ac{RL} algorithms~\cite{Makoviychuk2021IsaacGym,isaac_sim_510_physics_fundamentals}.

\subsection{Suitability of Isaac Sim~5.1 and IsaacLab for cloth manipulation in RL}
\subsubsection{Integrated cloth and deformable simulation support}
For cloth manipulation, contact-rich interactions between articulated robots and deformable objects must be simulated with sufficient stability while maintaining the throughput required by \ac{RL}.
PhysX provides \ac{GPU}-centric simulation features beyond rigid-body dynamics, including particles and soft-body related capabilities; in the \ac{GPU} simulation documentation, it is explicitly stated that soft bodies and particles are \ac{GPU}-only features~\cite{physx_541_gpu_rigid_bodies}.
In addition, the same documentation states that a position-based particle system can be used to simulate cloth and other deformable entities, enabling direct interaction with rigid bodies and articulated mechanisms within a unified physics scene~\cite{physx_541_gpu_rigid_bodies}.
Within Isaac Sim, physics properties are encoded at the \ac{USD} level via schemas and are executed by the PhysX backend, which allows deformable assets and robot assets to be configured in a consistent authoring model and trained in a shared simulation stage~\cite{isaac_sim_510_physics_fundamentals,isaac_sim_510_overview}.
This integration is aligned with the thesis objective of training manipulation policies that interact with cloth while also requiring accurate robot articulation dynamics and contact handling.

\subsubsection{Scalable training and distributed execution}
High-throughput training requires that many environments be advanced in parallel, typically under headless execution with minimized rendering overhead.
Isaac Sim explicitly supports accelerated simulation without rendering, which is essential when the observation model is state-based and the training objective is dominated by physics stepping rather than image synthesis~\cite{isaac_sim_510_physics_fundamentals}.
On top of Isaac Sim, IsaacLab provides multi-\ac{GPU} and multi-node execution paths to scale training workloads, including orchestration patterns for launching multiple processes and assigning environments to devices~\cite{isaaclab_multi_gpu_doc}.
Such scaling capabilities are practically relevant when the effective sample throughput of a single device becomes the limiting factor, as has been observed in large-scale robot learning contexts that rely on massive parallel simulation to collect sufficient experience for deep \ac{RL}~\cite{Makoviychuk2021IsaacGym}.

\subsubsection{Implications for the present thesis}
For the present thesis, Isaac Sim~5.1 and IsaacLab jointly address three constraints that arise from cloth manipulation in \ac{RL} settings:
\begin{itemize}
    \item \textbf{Unified rigid--deformable contact simulation:} A single PhysX-backed simulation scene can contain articulated robots and deformable cloth objects, reducing integration friction between robot modelling and cloth dynamics configuration~\cite{isaac_sim_510_physics_fundamentals,physx_541_gpu_rigid_bodies}.
    \item \textbf{Throughput-oriented execution:} Headless and accelerated simulation modes enable training loops where physics stepping dominates runtime, which is desirable for state-based policies and large-scale data collection~\cite{isaac_sim_510_physics_fundamentals,Makoviychuk2021IsaacGym}.
    \item \textbf{Reusable task formalization:} The manager-based \ac{MDP} workflow provides a structured decomposition of task components (observations, rewards, terminations, events), which supports iterative task development and ablation-style experimentation~\cite{isaaclab_mdp_api}.
\end{itemize}

\subsection{Comparison with alternative simulation platforms}
The choice of simulator impacts both physical fidelity and the feasibility of large-scale training.
A qualitative comparison with commonly used alternatives is summarized in Table~\ref{tab:simulator_comparison}, with emphasis on properties that are critical for cloth manipulation \ac{RL} workloads.

MuJoCo is widely used in robotics and control research and is designed as a physics engine for model-based control, providing efficient simulation of articulated rigid-body systems~\cite{Todorov2012MuJoCo}.
Its standard usage emphasizes fast and stable rigid-body dynamics, which makes it attractive for control-centric benchmarks, particularly when \ac{MJCF}-based models are readily available~\cite{Todorov2012MuJoCo}.
However, the present thesis requires integrated deformable dynamics for cloth in the same training loop, which motivates selecting a simulator stack that exposes deformable/particle-based mechanisms in conjunction with articulated robot models~\cite{physx_541_gpu_rigid_bodies}.

PyBullet is a Python interface to the Bullet physics engine and is commonly used for robotics prototyping and learning research due to its scripting convenience and support for standard robot description formats~\cite{pybullet_home}.
In typical usage, PyBullet enables rapid iteration on rigid-body manipulation tasks, but large-scale throughput comparable to \ac{GPU}-accelerated pipelines is not the primary design target of this ecosystem~\cite{pybullet_home}. % TODO: find citation on (lack of) GPU-parallel batched simulation focus

SOFA is an open-source framework that has been developed with a strong emphasis on physically detailed simulation for medical and deformable-object scenarios, including a modular architecture for combining multiple physical models and solvers~\cite{Allard2007SOFA}.
This focus makes SOFA a natural candidate when high-fidelity deformable simulation is prioritized.
For the learning-centric setting considered here, the practical bottleneck is the combination of deformable dynamics with high-throughput data collection, where a tight integration with \ac{GPU}-accelerated simulation and batched environment execution is typically required~\cite{Makoviychuk2021IsaacGym}. % TODO: find citation comparing SOFA throughput vs. GPU-parallel RL simulators

In summary, Isaac Sim and IsaacLab are selected because (i) a single solver stack is provided for articulated rigid bodies and cloth-like deformables, (ii) \ac{GPU}-accelerated execution paths are exposed for throughput-oriented training, and (iii) a reusable \ac{MDP} workflow is supported for systematic task development and distributed training~\cite{isaac_sim_510_physics_fundamentals,physx_541_gpu_rigid_bodies,isaaclab_mdp_api,isaaclab_multi_gpu_doc}.

\begin{table}[htb]
    \centering
    \caption{Qualitative comparison of simulation platforms for learning-based cloth manipulation.}\label{tab:simulator_comparison}
    \begin{tabular}{l c c c}
    \hline
    \textbf{Criterion} & \textbf{Isaac Sim + IsaacLab} & \textbf{MuJoCo} & \textbf{PyBullet / SOFA} \\
    \hline
    Articulated rigid-body control & yes~\cite{isaac_sim_510_physics_fundamentals} & yes~\cite{Todorov2012MuJoCo} & yes~\cite{pybullet_home} \\
    Deformable/cloth dynamics in same stack & yes~\cite{physx_541_gpu_rigid_bodies} & limited % TODO: verify claim
    & yes (SOFA focus)~\cite{Allard2007SOFA} \\
    GPU-parallel simulation for RL throughput & supported~\cite{physx_541_gpu_rigid_bodies,isaaclab_multi_gpu_doc} & not primary focus~cite{Todorov2012MuJoCo}% TODO: find citation
    & not primary focus~\cite{pybullet_home} % TODO: find stronger citation \\
    Photorealistic RTX rendering & yes~\cite{isaac_sim_510_overview,omniverse_rtx_renderer_doc} & no % TODO: find citation
    & limited % TODO: verify claim \\
    ROS~2 integration workflows & yes~\cite{isaac_sim_ros2_reference_architecture_510} & external wrappers % TODO: find citation
    & external wrappers % TODO: find citation \\
    \hline
    \end{tabular}
\end{table}

% ---------------------------------------------------------------------------
% BibTeX keys used in this section:
%  - isaac_sim_510_overview                 (NVIDIA Isaac Sim 5.1 overview: Omniverse-based app, USD unification, PhysX + RTX, links to Isaac Lab)
%  - isaac_sim_510_release_notes            (NVIDIA Isaac Sim 5.1 release notes)
%  - isaac_sim_510_physics_fundamentals     (NVIDIA Isaac Sim physics fundamentals: USD Physics/PhysX schemas and simulation execution)
%  - isaac_sim_ros2_reference_architecture_510 (NVIDIA Isaac Sim ROS 2 reference architecture and bridge-based integration)
%  - omniverse_rtx_renderer_doc             (Omniverse RTX Renderer documentation)
%  - physx_541_gpu_rigid_bodies             (PhysX 5.4.1 documentation: GPU rigid bodies and GPU-only features including soft bodies/particles)
%  - isaaclab_overview_doc                  (Isaac Lab documentation overview and project positioning)
%  - isaaclab_mdp_api                       (Isaac Lab API documentation for MDP and manager responsibilities)
%  - isaaclab_multi_gpu_doc                 (Isaac Lab documentation for multi-GPU and multi-node training)
%  - Mittal2023Orbit                        (Orbit framework paper in IEEE RA-L)
%  - Makoviychuk2021IsaacGym                (Isaac Gym paper: GPU-accelerated simulation for robot learning)
%  - Todorov2012MuJoCo                      (MuJoCo paper: physics engine for model-based control)
%  - pybullet_home                          (PyBullet official project website)
%  - Allard2007SOFA                         (SOFA framework paper: open-source framework for medical simulation)
% ---------------------------------------------------------------------------

% ---------------------------------------------------------------------------
% BibTeX entries — add to bibliography/literature.bib:
%
% @online{isaac_sim_510_overview,
%   author  = {{NVIDIA}},
%   title   = {What Is Isaac Sim?},
%   year    = {2026},
%   url     = {https://docs.isaacsim.omniverse.nvidia.com/5.1.0/index.html},
%   urldate = {2026-03-03},
% }
%
% @online{isaac_sim_510_release_notes,
%   author  = {{NVIDIA}},
%   title   = {Release Notes},
%   year    = {2026},
%   url     = {https://docs.isaacsim.omniverse.nvidia.com/5.1.0/release_notes.html},
%   urldate = {2026-03-03},
% }
%
% @online{isaac_sim_510_physics_fundamentals,
%   author  = {{NVIDIA}},
%   title   = {Physics Simulation Fundamentals},
%   year    = {2026},
%   url     = {https://docs.isaacsim.omniverse.nvidia.com/5.1.0/physics/simulation_fundamentals.html},
%   urldate = {2026-03-03},
% }
%
% @online{isaac_sim_ros2_reference_architecture_510,
%   author  = {{NVIDIA}},
%   title   = {ROS 2 Reference Architecture},
%   year    = {2026},
%   url     = {https://docs.isaacsim.omniverse.nvidia.com/5.1.0/ros2_tutorials/reference_architecture.html},
%   urldate = {2026-03-03},
% }
%
% @online{omniverse_rtx_renderer_doc,
%   author  = {{NVIDIA}},
%   title   = {Omniverse RTX Renderer},
%   year    = {2026},
%   url     = {https://docs.omniverse.nvidia.com/materials-and-rendering/latest/rtx-renderer.html},
%   urldate = {2026-03-03},
% }
%
% @online{physx_541_gpu_rigid_bodies,
%   author  = {{NVIDIA}},
%   title   = {GPU Simulation (PhysX 5.4.1 Documentation)},
%   year    = {2026},
%   url     = {https://nvidia-omniverse.github.io/PhysX/physx/5.4.1/docs/GPURigidBodies.html},
%   urldate = {2026-03-03},
% }
%
% @online{isaaclab_overview_doc,
%   author  = {{The Isaac Lab Project Developers}},
%   title   = {Overview},
%   year    = {2026},
%   url     = {https://isaac-sim.github.io/IsaacLab/v1.0.0/index.html},
%   urldate = {2026-03-03},
% }
%
% @online{isaaclab_mdp_api,
%   author  = {{The Isaac Lab Project Developers}},
%   title   = {isaaclab.envs.mdp (API Documentation)},
%   year    = {2026},
%   url     = {https://isaac-sim.github.io/IsaacLab/main/source/api/lab/isaaclab.envs.mdp.html},
%   urldate = {2026-03-03},
% }
%
% @online{isaaclab_multi_gpu_doc,
%   author  = {{The Isaac Lab Project Developers}},
%   title   = {Multi-GPU and Multi-Node Training},
%   year    = {2026},
%   url     = {https://isaac-sim.github.io/IsaacLab/main/source/features/multi_gpu.html},
%   urldate = {2026-03-03},
% }
%
% @article{Mittal2023Orbit,
%   author  = {Mittal, Mayank and Yu, Felix and Zhao, Jitendra and Cheng, Yuanhao and Dimitrov, Dimitar and Qiu, Louis and Velonaki, Mari and Zhu, Yuke},
%   title   = {Orbit: A Unified Simulation Framework for Interactive Robot Learning Environments},
%   journal = {IEEE Robotics and Automation Letters},
%   volume  = {8},
%   number  = {6},
%   pages   = {3740--3747},
%   year    = {2023},
%   doi     = {10.1109/LRA.2023.3270034},
% }
%
% @inproceedings{Makoviychuk2021IsaacGym,
%   author    = {Makoviychuk, Viktor and Wawrzyniak, Lukasz and Guo, Yunrong and Lu, Michelle and Storey, Kyle and Macklin, Miles and Hoeller, David and Rudin, Nikita and Allshire, Arthur and Handa, Ankur and State, Gavriel and Girdhar, Ankur and Kumar, Tushar and Fink, Jonathan and Merel, Josh and Tassa, Yuval},
%   title     = {Isaac Gym: High Performance GPU Based Physics Simulation for Robot Learning},
%   booktitle = {NeurIPS 2021 Datasets and Benchmarks Track (Round 2)},
%   year      = {2021},
%   url       = {https://openreview.net/forum?id=fgFBtYgJQX_},
% }
%
% @inproceedings{Todorov2012MuJoCo,
%   author    = {Todorov, Emanuel and Erez, Tom and Tassa, Yuval},
%   title     = {MuJoCo: A Physics Engine for Model-Based Control},
%   booktitle = {2012 IEEE/RSJ International Conference on Intelligent Robots and Systems (IROS)},
%   pages     = {5026--5033},
%   year      = {2012},
%   doi       = {10.1109/IROS.2012.6386109},
% }
%
% @online{pybullet_home,
%   author  = {{Bullet Physics Project}},
%   title   = {Bullet Real-Time Physics Simulation: PyBullet},
%   year    = {2026},
%   url     = {https://pybullet.org/wordpress/},
%   urldate = {2026-03-03},
% }
%
% @article{Allard2007SOFA,
%   author  = {Allard, J. and Cotin, S. and Faure, F. and Bensoussan, P.-J. and Poyer, F. and Duriez, C. and Delingette, H. and Grisoni, L.},
%   title   = {SOFA -- an open source framework for medical simulation},
%   journal = {Studies in Health Technology and Informatics},
%   volume  = {125},
%   pages   = {13--18},
%   year    = {2007},
%   pmid    = {17377224},
% }
% ---------------------------------------------------------------------------